%%%%%%%%%%%%%%%%%%%%%%%%%%%%%%%%%%%%%%%%%%%%%%%%%%%%%%%%%%%
% --------------------------------------------------------
% FAPET UGM Rho
% LaTeX Template
% Version 1.0.0 (18/10/2025)
%
% Adapted for:
% Fakultas Peternakan UGM
% --------------------------------------------------------
%%%%%%%%%%%%%%%%%%%%%%%%%%%%%%%%%%%%%%%%%%%%%%%%%%%%%%%%%%%

\documentclass[9pt,article,twoside]{fapet-rho-class/fapet-rho}
\FAPETtemplatetype{\FAPET} % Select your article type

%----------------------------------------------------------
% JOURNAL/HEADER INFORMATION
%----------------------------------------------------------

\vol{49}
\issue{3}
\setcounter{page}{1}
\pages{1-8}
\monthyear{October 2025}
\thisyear{2025}
\doi{\href{http://buletinpeternakan.fapet.ugm.ac.id/}{http://buletinpeternakan.fapet.ugm.ac.id/}}

%----------------------------------------------------------
% TITLE
%----------------------------------------------------------

\title{Reproductive Performance of Madura Cows at Various Grades in the Breeding Center of Pamekasan Regency}

%----------------------------------------------------------
% AUTHORS AND AFFILIATIONS
%----------------------------------------------------------

\author[1,$*$]{Faigah Musdalifah}
\author[1]{Mudawamah}
\author[1]{Nurul Humaidah}

%----------------------------------------------------------

\affil[1]{Magister program of Animal science, Department of Animal Science, University of Islam Malang, Malang 65144, Indonesia}

%----------------------------------------------------------
% CORRESPONDING AUTHOR INFORMATION
%----------------------------------------------------------

\leadauthor{Musdalifah et al.}
\smalltitle{Reproductive Performance of Madura Cows}

%----------------------------------------------------------
% ARTICLE INFORMATION
%----------------------------------------------------------

\corres{Faigah Musdalifah}
\email{faigahmusdalifah02@gmail.com; mudawamah@unisma.ac.id}

\received{September 15th, 2025}
\accepted{October 18th, 2025}

\license{This work is licensed under a Creative Commons Attribution 4.0 International License}

\setbool{rho-abstract}{true} % Set false to hide the abstract

%----------------------------------------------------------
% ABSTRACT AND KEYWORDS
%----------------------------------------------------------

\begin{abstract}
The purpose of this study was to analyze the reproductive performance of Madura cattle at various grades in the breeding center of Pamekasan Regency. The method used was a case study with purposive sampling, with the criteria being 120 Madura cattle grouped by grade with an age range of 30-36 months. Reproductive performance data were analyzed using the Chi-square test. The reproductive performance of Madura cows at various grades compared to the ideal reproductive performance standard was not significantly different (P>0.05). The average reproductive performance of grade 1, 2, and 3 Madura cattle at the first estrus age (17.30±1.02 months, 17.48±1.04 months, and 17.96±1.13 months), first mating age (20.17±1.74 months, 20.52±1.89 months, and 20.90±2.04 months), first calving age (30.43±2.25 months, 30.76±2.26 months, and 30.98±2.18 months), S/C value (1.30±0.60 times, 1.33±0.61 times, and 1.40±0.61 times) and CR value (76.67\%, 73.81\%, and 66.67\%). The reproductive performance of Madura cattle at various grades still meets the ideal standard, but there is a tendency for the average reproductive performance of Madura cattle from the study to be better than the ideal standard.
\end{abstract}

\keywords{cow, reproductive performance, Madura cattle}

%----------------------------------------------------------

\begin{document}

\maketitle
\pagestyle{fancy}\thispagestyle{firststyle}

%----------------------------------------------------------

\section{INTRODUCTION}

Madura cattle are one of the breeds of cattle that are widely integrated into the socio-cultural and socio-economic life of society, for example for traction, savings, or as a means of sport and entertainment such as kerapan and sonok cattle. Madura cattle have special value for the Madurese people for their social status and have considerable potential as a tourist attraction \citep{Zali2019}. The selection of Madura cattle, particularly in the Madura cattle Breeding Center of Pamekasan Regency, involves dividing cattle into groups with specific qualifications known as grades. Grade determination is based on the Pamekasan Regency Madura Cattle Regional Standard. Grade identification is a grouping based on cattle performance \citep{Putra2018}. The challenge in maintaining local cattle is maintaining the population of quality Madura cattle to prevent negative selection, which results in the exclusion of good-quality cattle, leaving only low-quality or poor-quality cattle \citep{Susilawati2017}. Efforts to improve selection methods are essential to increase local cattle productivity. One method of selecting breeding stock is by observing the livestock's reproductive performance.

Efforts to improve the genetic quality of livestock can be carried out by increasing the efficiency of livestock reproduction. Reproductive performance can indirectly affect livestock productivity \citep{Pariswara2021}. Livestock reproductive information is used to determine and identify reproductive capabilities that can support population increase efforts \citep{Makmur2019} and as a determinant of livestock productivity \citep{Sartika2020}. Reproductive performance is a parameter that can indicate whether a livestock's reproductive activity is running well or not. In this case, it can include age at first estrus, age at first mating, age at first calving, Service per Conception (S/C), and Conception Rate (CR). Reproductive studies on Madura cattle have generally been conducted without distinguishing or classifying based on morphometric grades. This study aims to analyze the reproductive performance of Madura cattle according to morphometric grades, in order to provide a more specific overview of the reproductive potential across at various grades in the breeding center of Pamekasan Regency.

\section{MATERIALS AND METHODS}

The research was conducted in the Breeding Center in Pasean District, Pamekasan Regency from February to April 2025. The method used in this research is a case study. The data collected in this study comprised both primary and secondary data. Primary data was obtained from interviews and morphometric measurements. The study involved only non-invasive morphometric measurements and observational data collection. All procedures complied with animal welfare standards, and no undue stress or harm was inflicted on the cattle. Secondary data was obtained from relevant agency documents and literature or references relevant to the research. The variables used in this study are reproductive performance which includes first estrus age, first mating age, first calving age, Service per Conception (S/C) and Conception Rate (CR). The population in this study were Madura cattle in the breeding center, Pasean District, Pamekasan Regency. The sample was determined based on Madura cattle population data from Dinas Ketahanan Pangan dan Pertanian Pamekasan Regency. The sample was taken using purposive sampling, with the criteria being Madura cattle grouped by grade with an age range of 30-36 months. The sample used was 120 Madura cattle, with grade determination in accordance with the 2022 Pamekasan Regency Madura Cattle Regional Standard concerning the minimum quantitative requirements for Madura cattle breed.

This study used several instruments to collect the necessary data. The instruments used included measuring instruments and observation sheets. The measuring instruments used in the study were a measuring tape and a measuring stick, which were used to measure livestock body size. The measuring stick was used to measure body length and shoulder height, while the measuring tape was used to measure chest circumference. The observation sheet was used to record the results of the research data obtained. Grade identification begins with individual body size data obtained and then compared with regional standards as listed in Table 1. The basis for grade assessment is the minimum value of the standard that must be achieved. An example of how to identify grades is presented in Table 2.

Service per Conception (S/C) and Conception Rate (CR) formula for Madura cattle includes:

\begin{enumerate}
    \item Service per Conception (S/C) is the number of matings carried out to produce one pregnancy which can be calculated using the following formula:
    
    \begin{equation}
    S/C = \frac{\text{Number of cattle mated until pregnancy occurs}}{\text{Number of pregnant cattle}}
    \end{equation}
    
    \item Conception Rate (CR) or the percentage of cattle that become pregnant during the first mating or artificial insemination. CR calculation formula can be formulated as follows:
    
    \begin{equation}
    CR = \frac{\text{Number of pregnant cattle at first mating}}{\text{Number acceptors}} \times 100\%
    \end{equation}
\end{enumerate}

Reproductive performance data of Madura cattle at various grades were analyzed using the Chi-square test to compare research results and ideal standards with the following formula:

\begin{equation}
\chi^2 = \sum \frac{(O-E)^2}{E}
\end{equation}

\noindent Description:

\begin{tabular}{ll}
$\chi^2$ & = Chi-square test statistics \\
$O$ & = Observed frequency \\
$E$ & = Expected frequency \\
\end{tabular}

\section{RESULTS AND DISCUSSION}

\subsection{First Estrus Age, First Mating Age and First calving age for Madura Cattle at Various Grades}

The results of the research on first estrus age, first mating age and first calving age of Madura cattle at various grades can be seen in Table 3. The results of Chi-square statistical analysis test showed that first estrus age of Madura cattle at various grades with ideal standard of first estrus of Madura cattle was not significantly different (P>0.05). This shows that first estrus age of Madura cattle is still in accordance with ideal standard of first estrus of Madura cattle, but average first estrus age of Madura cattle grades 1, 2 and 3 of the study had a tendency to be better by 3.89\%, 2.91\% and 0.23\% compared to ideal standard or shorter by about 21 days, 15 days, and one day. The first estrus age of Madura cattle grades 1, 2 and 3 were 17.30±1.02 months, 17.48±1.04 months, and 17.96±1.13 months, respectively. The research results are lower than the Decree of the Minister of Agriculture (2010) regarding determination of Madura cattle breed which reports that puberty age of Madura cattle is 549-610 days or around 18-20 months.

One factor thought to influence age at first estrus is BCS or Body Condition Score. Cattle with an ideal BCS tend to experience their first estrus at an ideal age. This is consistent with the opinion of \citet{Heryani2019} who stated that first estrus age indicates a good BCS. \citet{Wicaksana2024} added that first estrus age is influenced by several factors, including the cow's nutrition from weaning to heifer, genetics, and husbandry management. Furthermore, differences in cattle breeds also influence how quickly or slowly a cow experiences its first estrus.

The results of the Chi-square statistical analysis test showed that first mating age of Madura cattle at various grades with ideal standard of first mating age of female cattle was not significantly different (P>0.05). This shows that first mating age of Madura cattle is still in accordance with ideal standard of first mating age of Madura cattle, but there is a tendency for average first mating age of Madura cattle grades 1, 2 and 3 to be better by 8.33\%, 6.71\% and 5.02\% compared to ideal standard or shorter by around 55 days, 44 days and 33 days. Female cattle are expected to start mating for first time at age of between 18-22 months \citep{Ismaya2014}. The results of the study on first mating age in Madura cattle grades 1, 2 and 3 were 20.17±1.74 months, 20.52±1.89 months, and 20.90±2.04 months, respectively. This result is lower than \citet{Kutsiyah2017}, which stated that average first mating age for female Madurese cattle was 21.91±5.73 months. \citet{Wisono2015} also reported that average first mating age and first calving for Madura cattle in Tlanakan, Pegantenan, and Pasean Districts were 23.79±2.53 months, 23.06±2.86 months, and 23.48±2.31 months, respectively.

First mating age in livestock is closely related to adequate physical maturity, thus providing a greater opportunity to produce a large number of offspring. At the research site, most farmers did not immediately mate their cattle at first estrus and tended to wait several more estruses before mating. This is because farmers wait for their cattle to reach optimal physical maturity and reproductive organ maturity so that the female's body is ready for pregnancy and parturition. According to \citet{Wicaksana2024}, one factor causing abnormal births is cows that are mated too young.

Selecting right timing for breeding heifers will impact production cost efficiency by reducing expenses and increasing reproductive effectiveness \citep{Wicaksana2024}. Age and body size of heifers at first mating need to be considered to achieve a high conception rate and prevent difficulties during calving. If heifers are not mated at ideal age and body size, their reproductive performance tends to be low \citep{Desinawati2010}.

Results of the Chi-square statistical analysis test showed that first calving age of Madura cattle at various grades with ideal standard of first calving of Madura cattle was not significantly different (P>0.05). This shows that first calving age of Madura cattle is still in accordance with ideal standard of first calving of Madura cattle, but there is a tendency for the average first calving age of Madura cattle grades 1, 2 and 3 to be better by 1.83\%, 0.77\% and 0.07\% compared to the ideal standard or shorter by about 17 days, 7 days and less than a day. According to \citet{Ismaya2014} that in beef cattle, it is expected to give birth to the first calf at age of 27-31 months. The results of the first calving age in Madura cattle grades 1, 2 and 3 were 30.43±2.25 months, 30.76±2.26 months, and 30.98±2.18 months, respectively. On average, farmers in the study area mate their cattle at less than two years of age. This result is lower than \citet{Kutsiyah2017}, which stated that first calving age for Madura cattle was 31.97±6.43 months. \citet{Wisono2015} study also reported that average first calving age for Madura cattle in Tlanakan, Pegantenan, and Pasean Districts was 34.63±2.46 months, 33.44±2.95 months, and 33.74±2.65 months, respectively.

First calving age is influenced by several factors, including first mating age, S/C value, gestation length, and husbandry management. According to \citet{Pelmus2024}, first calving age is economically important. Genetic, environmental, and physiological factors play a significant role in increasing first calving age. Higher reproductive efficiency can increase meat production and livestock profitability. \citet{Awan2016} added that ideal first calving age can be achieved by providing adequate, high-quality feed from the calf to heifer stage. Insufficient feed or low nutrient content in feed can inhibit growth, thus prolonging age at first calving.

Implementing appropriate husbandry management during the growth period of young cattle or heifers directly influences growth rate and optimal body weight attainment. Achieving maximum body weight in a relatively short time will accelerate puberty and shorten time to first calving \citep{Sumiyanti2023}. \citet{Napitupulu2024} reported that first calving age in cattle is influenced by several factors, namely age at puberty, age at first mating, and length of gestation. \citet{Hartatik2009} also added that a high S/C value will result in a longer birth interval compared to normal conditions.

\subsection{Service Per Conception (S/C) and Conception Rate (CR) of Madura Cattle at Various Grades}

The results of S/C and CR research on Madura cattle at various grades are shown in Table 4. The results of statistical analysis showed that S/C value of Madura cattle at various grades compared to the ideal standard S/C of female cattle was not significantly different (P>0.05). Results of field observations obtained S/C value of Madura cattle at grades 1, 2, and 3, respectively, 1.30±0.60 times, 1.33±0.61 times, and 1.40±0.61 times. According to research by \citet{Kutsiyah2018}, S/C value of Madura cattle reached 1.40 times, while research by \citet{Wulansari2023} that S/C of Madura cattle was 1.5 times. The S/C value obtained in the study was still within ideal standard range of S/C, this indicates that S/C value in the study was good.

The results of S/C value of Madura cattle at various grades is still in accordance with ideal standard, but there is a tendency for average S/C value of Madura cattle grades 1, 2 and 3 research results to be better by 18.75\%, 16.67\% and 12.76\% compared to ideal standard S/C of female cattle. According to \citet{Tophianong2023}, ideal S/C value is between 1.6-2. \citet{Khan2015}, stated that S/C is number of natural or artificial mating services required until pregnancy occurs. S/C value is a simple criterion for determining fertility and analyzing costs. Livestock with 1-2 services showed a CR of 62.4\% and livestock that required 3 or more times showed a CR of 50\%.

One factor that influences S/C value is the Body Condition Score (BCS). BCS plays a crucial role in determining a cow's reproductive efficiency. An ideal BCS increases the chance of pregnancy, thus lowering the S/C rate. According to \citet{Harmayani2023}, one factor that influences reproductive efficiency in cows is BCS value. Cows with a BCS of 3-4 are included in ideal BCS category because they are neither too fat nor too thin. \citet{Putra2023} added that a low BCS due to feed deficiency can cause a delay in the onset of estrus and prolong the duration of estrus, thus impacting an increase in S/C value.

Good nutritional intake and adequate feed intake can influence reproductive process in livestock. Furthermore, ideal S/C value is also influenced by knowledge possessed by farmers. Several farmers in breeding center in Pasean District can already detect estrus by observing signs shown by cows such as restlessness, decreased appetite, swollen, slimy, and red vaginas. The importance of farmer knowledge is closely related to appropriate mating timing, thus increasing chances of pregnancy. The S/C value is influenced by several factors, including feed, farmer knowledge, semen quality, and the health of dam. The S/C value indicates the fertility of the livestock, with a high S/C value indicating a lower level of fertility \citep{Wanma2022}. According to \citet{Jurame2018}, failure to detect estrus can result in pregnancy failure. Therefore, accurate estrus detection is crucial in mating programs to ensure timely fertilization.

The results of the Chi-square statistical analysis showed that CR values of Madura cattle at various grades with the ideal standard CR of female cattle were not significantly different (P>0.05). This indicates that CR values of Madura cattle at various grades are still in accordance with ideal standard, but there is a tendency for percentage of CR values of Madura cattle grades 1, 2 and 3 to be better than ideal standard CR of female cattle, 11.67\%, 8.81\% and 1.67\%, respectively. According to \citet{Tophianong2023}, ideal standard CR value ranges from 65-75\%. The results of field observations obtained CR values of Madura cattle at grades 1, 2, and 3 were 76.67\%, 73.81\%, and 66.67\%, respectively. Value of CR obtained in this study are in accordance with ideal standard CR value. \citet{Wulansari2023}, CR of Madura cattle is 66.667\%.

The results of CR value of Madura cattle of various grades in study area is thought to be influenced by nutritional factors, livestock health, farmer knowledge in detecting estrus, and proper mating management. Furthermore, an ideal BCS supports hormonal balance and increases chance of pregnancy, thereby increasing the CR percentage in cattle. Percentage of CR is influenced by several factors such as nutritional intake and fertility factor of livestock \citep{Prasdini2018}. Furthermore, CR value can also be influenced by Body Condition Score (BCS) and timing of mating. Cows with a BCS ≥ 3 show a higher CR value \citep{Martini2022}. \citet{Samkange2019} added that seasonal factors influence CR value, with winter or rainy seasons thought to increase CR value. Value of CR is determined from results of livestock pregnancy diagnosis as seen from absence of signs of estrus \citep{Asfar2023}. Pregnancy in livestock is determined by rectal palpation of genital tract on days 60–90 post-mating. One factor influencing CR is health status of the female's reproductive organs \citep{Khan2015}.

\section{CONCLUSION}

The reproductive performance of Madura cattle at various grades is still in accordance with the ideal standard, but there is a tendency for the average reproductive performance of Madura cattle from the research results to be better than the ideal standard. Madura cattle grade 1 have better reproductive performance than Madura cattle grade 2 and 3.

%----------------------------------------------------------
% CONFLICT OF INTEREST
%----------------------------------------------------------

\section*{Conflict of Interest}

The researcher declared that they have no conflict of interest related to this study.

%----------------------------------------------------------
% FUNDING STATEMENT
%----------------------------------------------------------

\section*{Funding Statement}

This research received no specific grant from any funding agency in the public, commercial, or not-for-profit sectors.

%----------------------------------------------------------
% ACKNOWLEDGMENTS
%----------------------------------------------------------

\section*{Acknowledgments}

The authors wish to thank everyone who provided insight and expertise that greatly assisted the research.

%----------------------------------------------------------
% AUTHOR'S CONTRIBUTION
%----------------------------------------------------------

\section*{Author's Contribution}

All authors contributed equally to the conception, design, data collection, analysis, and writing of this manuscript.

%----------------------------------------------------------
% ETHICS APPROVAL
%----------------------------------------------------------

\section*{Ethics Approval}

The study involved only non-invasive morphometric measurements and observational data collection. All procedures were conducted in accordance with animal welfare principles and complied with Indonesian Law No. 18 of 2009 on Animal Husbandry and Animal Health, as amended by Law No. 41 of 2014. All procedures complied with animal welfare standards, and no undue stress or harm was inflicted on the cattle.

%----------------------------------------------------------
% REFERENCES
%----------------------------------------------------------

\bibliographystyle{fapet-rho-class/fapet}
\bibliography{madura}

%----------------------------------------------------------

\end{document}
